

This paper demonstrates how replacing grading discretion from a central examiner to the the local educators of students reduces the information content of grades. The new grades reflected the quality of the grade rewarding institution, with lower quality school adding higher grades to than higher quality schools. The grading standards rapidly evolves upwards with educational standards and the information content deteriorating in an increasing rate, with profit driven schools among the schools with fastest adjustments to the new grading regime. The loosening of grading standards are also accompanied by a pattern for teachers/schools to reward grades that satisfies the acceptance requirement provided by the universities. The grade matching adds to another source of information deterioration of the rewarded grade, as more institutions with stricter grade requirement are subject to a worse deterioration of grade standards as schools/teachers adjust for the universities responses. 

Despite the loosening of grading standard was larger in lower quality institutions the impact of the jump in grade were limited to placement improvements for students from disadvantaged socio-economic background. Additionally the improvements in representations at selective institutions were limited to students above the average in the pre-pandemic grade distribution. The non-monotone relationship between grade and placement mobility were driven by differences in application choices between students with varying baseline backgrounds, as disadvantaged students do not apply to selective institutions despite their rewarded grades facilitating their placements into selective institutions than students from more advantaged backgrounds. 

My results points to the dangers of accommodating test-optional policies into university admission systems. Proponents of test-optional policies advocate for a holistic admission system to diversify students background at selective universities, as well as using schools and teachers as a method for extracting information on students quality and potential that would be missed out in traditional standardized exams in which the examiner has no-knowledge of the students background when they are grading the tests. The grading patterns in the UK demonstrates the vulnerability of such decentralized grading regimes to potential gaming of grades from teachers, as well as their limited effect on diversifying selective universities with students socio-economic backgrounds. The speed of the adaptation were fastest for private schools, which teaches the students from the wealthiest background of parents. I suggest that simply lowering the barriers to selective universities will not improve the diversity of selective universities. Even among the students that apply to universities, students from socio-economics advantaged vary with their attitudes towards selective institutions, such as the composition of students or their parents attitudes or knowledge on elite institutions. I propose that the lowering the barriers for universities to diversity their knowledge to students from diverse backgrounds could be met by reducing frictions in the benefits of higher-education but not on the application system.

A important caveat of the policy implication is whether policy makers can take advantage of the regressive grade inflation from non-standardized testing while assuring that the increase in grades for socio-economically disadvantaged students are converted to improvements into placements into selective institutions. As my results uses the changes in grades due to COVID19 pandemic, the application patterns that I use for assessing placement pattern may be influenced by the non-normal situation accruing from the pandemic chaos. For example, students from disadvantaged backgrounds may refrain from applying ambitiously despite receiving high grades, and also students may prefer to study at universities closer to their parents to ease the civil unrest during the pandemic. A feasible extension of my results would be to take advantage of the growing literature on estimating of university application portfolio choice (\cite{10.1257/aer.20211137}) and test whether placement patterns would change if students had changed their application choice by foreseeing the changes in the grading scheme.

